Na vele maanden varen gingen ze aan land.
Emmy was inderdaad in Marrakesh gestrand.

De haven lag direct aan een marktplein.
Blijkbaar moest ook zij daar zijn.

Het ruim werd aldaar geleegd.
Ontucht werd overal gepleegd.

Onder het mom eerst voelen dan kopen,
werd er flink in de handelswaar gekropen.

De mastdames waren gratis te proberen.
Als lokketjes om de koopwaar aan te smeren.

Emmy werd op een apart podium geplaatst.
Veel aanwezigen verschenen gehaast

om op Emmy te kunnen bieden.
Er waren wel vijfhonderd kooplieden.

Toen de veilingmeester haar jurk af liet glijden
Wist dat de kopers nog meer te verblijden.

De kapitein was verheugd over de biedingen.
Het beloonde zijn lastige inspanningen

om niet het meisje tijdens de reis te nemen.
Dat zou haar alle waarde ontnemen.

Immers geen Berber wil een bevlekte  vrouw.
Hij spelde de menigte op de mouw

dat Emmy nog nooit wat aangetast.
Hij had haar geschiedenis aangepast.

Voor Moren was een roodharige een buitenaards wezen.
Van een Noordeling kon niet het maagdenvlies worden gelezen.

Zoals op de wijze bij de Afrikanen of Arabieren.
Met dit fabeltje wist men de handel te versieren.

Ontmaagde dellen lang onaangeraakt
werden in Europa overal geschaakt.

Vooral de stad van de jager was geliefd.
Daar werd de koopwaar grondig debrieft.

Als ze niet zeker waren over de droogstand
dan werd de sloerie in het klooster geplant.

Of men wist – zoals bij Emmy -
bezwangering was geen optie.

Men moest tijdens de reis van de handelswaar afblijven.
Anders groeide er misschien een kindje in hun lijven.

Daarmee zou de mythe worden ontkracht.
Daarop waren de kapiteins zeer bedacht.

De keurmeester zat in het complot.
Hij stak zijn vingers in Emmy’s grot

en knikte bevestigend op haar maagdelijkheid.
Ook al was Emmy die al twee jaar kwijt.

De koper was een handelaar
ergens in de tachtig jaar.

Emmy dacht; “wat heb ik met oude mannen.
Wil geen jongeman zich voor me inspannen?”
